


\documentclass[11pt, a4paper]{article}
\usepackage[slovenian]{babel}\usepackage{csquotes}
\usepackage[T1]{fontenc}



\usepackage{geometry}
%\setlength{\parindent}{0pt}
\usepackage{parskip}
\usepackage{tikz, graphicx, float}
\usepackage{enumerate}

\usepackage{hyperref}

\usepackage{titlesec}
\titleformat{\section}{\large}{\thesection.}{0.5em}{}%\centering}


\usepackage{amsmath, amsthm, amsfonts, amssymb, old-arrows, mathtools}

\theoremstyle{plain}{
    \newtheorem{izrek}{Izrek}
    \newtheorem{trditev}{Trditev}
    \newtheorem{lema}{Lema}
}
\theoremstyle{definition}{
    \newtheorem{definicija}{Definicija}
    \newtheorem{zgled}{Zgled}
}
\theoremstyle{remark}{
    \newtheorem{opomba}{Opomba}
}


\let\phi\varphi

\def\C{\mathbf{C}}
\def\F{\mathbf{F}}
\def\I{\mathbf{I}}
\def\R{\mathbf{R}}
\def\S{\mathbf{S}}
\def\T{\mathbf{T}}
\def\Z{\mathbf{Z}}


\usepackage{biblatex}
\addbibresource{main.bib}
\usepackage{url}


\title{Kako obesiti sliko}
\author{Maks Žnidaršič}
\date{}


\begin{document}

\maketitle


\section{Uvod}


Pokažemo pristop k konstrukciji \textit{prve homotopske grupe} skozi njeno kozmetično podobnost problemu obešanja slike (ang. picture hanging puzzle). Definicije in splošne izreke algebraične topologije povzamemo iz \cite{topologie-algebrique}, formulacija in rešitev problema pa sta predstavljeni v \cite{picture-hanging-puzzles}.

Problem si zastavimo takole: imamo sliko pritrjeno na sklenjeno vrvico, nekaj žebljev in prazno steno. Sliko z uporabo vrvice in žebljev pritrdimo na steno, kot kaže slika \ref{mucka}. \textit{Izziv je vrvico napeljati, tako da, ko iz stene odstranimo kateregakoli izmed žebljev, slika pade.}

\begin{figure}[H]
    \centering
    


\tikzset{every picture/.style={line width=.75pt}} %set default line width to 0.75pt        
\begin{tikzpicture}[x=.9pt,y=.9pt,yscale=-1,xscale=1]

\draw   (12.6,220.2) -- (172.6,220.2) -- (172.6,299.8) -- (12.6,299.8) -- cycle ;
\draw  [draw opacity=0] (52.35,279.49) .. controls (46.79,279.45) and (42.29,272.75) .. (42.29,264.49) .. controls (42.29,256.2) and (46.82,249.49) .. (52.4,249.49) .. controls (52.42,249.49) and (52.44,249.49) .. (52.45,249.49) -- (52.4,264.49) -- cycle ; \draw   (52.35,279.49) .. controls (46.79,279.45) and (42.29,272.75) .. (42.29,264.49) .. controls (42.29,256.2) and (46.82,249.49) .. (52.4,249.49) .. controls (52.42,249.49) and (52.44,249.49) .. (52.45,249.49) ;  
\draw  [draw opacity=0] (93.16,279.69) .. controls (98.37,279.31) and (102.51,272.8) .. (102.51,264.82) .. controls (102.51,256.63) and (98.14,249.98) .. (92.73,249.93) -- (92.68,264.82) -- cycle ; \draw   (93.16,279.69) .. controls (98.37,279.31) and (102.51,272.8) .. (102.51,264.82) .. controls (102.51,256.63) and (98.14,249.98) .. (92.73,249.93) ;  
\draw    (52.45,249.49) -- (59.18,233.27) ;
\draw    (59.18,233.27) -- (64.73,249.71) ;
\draw    (92.73,249.93) -- (87.18,234.16) ;
\draw    (87.18,234.16) -- (80.07,249.71) ;
\draw    (64.73,249.71) -- (80.07,249.71) ;
\draw  [draw opacity=0] (52.4,264.49) .. controls (52.58,261.63) and (54.72,259.38) .. (57.34,259.4) .. controls (60.05,259.41) and (62.24,261.87) .. (62.22,264.89) .. controls (62.22,264.9) and (62.22,264.92) .. (62.22,264.94) -- (57.3,264.86) -- cycle ; \draw   (52.4,264.49) .. controls (52.58,261.63) and (54.72,259.38) .. (57.34,259.4) .. controls (60.05,259.41) and (62.24,261.87) .. (62.22,264.89) .. controls (62.22,264.9) and (62.22,264.92) .. (62.22,264.94) ;  
\draw  [draw opacity=0] (82.86,264.37) .. controls (83.03,261.52) and (85.18,259.27) .. (87.8,259.28) .. controls (90.51,259.3) and (92.7,261.76) .. (92.68,264.77) .. controls (92.68,264.79) and (92.68,264.81) .. (92.68,264.82) -- (87.76,264.74) -- cycle ; \draw   (82.86,264.37) .. controls (83.03,261.52) and (85.18,259.27) .. (87.8,259.28) .. controls (90.51,259.3) and (92.7,261.76) .. (92.68,264.77) .. controls (92.68,264.79) and (92.68,264.81) .. (92.68,264.82) ;  
\draw  (58.91,240.56) -- (60.82,245.8) -- (57,245.8) -- cycle ;
\draw  (86.89,241.13) -- (88.78,246.6) -- (85,246.6) -- cycle ;
\draw    (38.51,268.38) -- (43.69,267.09) -- (46.4,266.42) ;
\draw    (106.29,269.49) -- (99.73,267.53) ;
\draw    (39.62,259.71) -- (46.18,261.53) ;
\draw    (99.18,262.16) -- (105.62,259.27) ;
\draw  [draw opacity=0] (95.69,250.94) .. controls (100.68,244.27) and (109.48,239.9) .. (119.45,240.02) .. controls (134.8,240.21) and (147.11,250.98) .. (146.95,264.08) .. controls (146.91,267.38) and (146.08,270.52) .. (144.61,273.36) -- (119.16,263.73) -- cycle ; \draw   (95.69,250.94) .. controls (100.68,244.27) and (109.48,239.9) .. (119.45,240.02) .. controls (134.8,240.21) and (147.11,250.98) .. (146.95,264.08) .. controls (146.91,267.38) and (146.08,270.52) .. (144.61,273.36) ;  
\draw  [draw opacity=0] (113.36,281.58) .. controls (108.99,280.18) and (106.23,278.13) .. (106.26,275.9) .. controls (106.32,271.88) and (115.37,268.75) .. (126.48,268.91) .. controls (134.52,269.03) and (141.43,270.84) .. (144.61,273.36) -- (126.37,276.19) -- cycle ; \draw   (113.36,281.58) .. controls (108.99,280.18) and (106.23,278.13) .. (106.26,275.9) .. controls (106.32,271.88) and (115.37,268.75) .. (126.48,268.91) .. controls (134.52,269.03) and (141.43,270.84) .. (144.61,273.36) ;  
\draw    (153,220.2) .. controls (151.8,186.2) and (197,180.6) .. (165,157) .. controls (133,133.4) and (126.6,199.8) .. (88.2,177) .. controls (49.8,154.2) and (0.41,167.93) .. (0.2,180.6) .. controls (-0.01,193.27) and (33.4,192.6) .. (33,219.8) ;
\draw  (68.61,266.2) .. controls (68.61,265.1) and (70.22,264.2) .. (72.21,264.2) .. controls (74.19,264.2) and (75.8,265.1) .. (75.8,266.2) .. controls (75.8,267.3) and (74.19,268.2) .. (72.21,268.2) .. controls (70.22,268.2) and (68.61,267.3) .. (68.61,266.2) -- cycle ;
\draw  [draw opacity=0] (49.46,281.23) .. controls (50.59,279.52) and (53.99,278.27) .. (58.02,278.25) .. controls (62.98,278.22) and (67.01,280.07) .. (67.02,282.38) .. controls (67.03,283.36) and (66.32,284.25) .. (65.13,284.97) -- (58.04,282.43) -- cycle ; \draw   (49.46,281.23) .. controls (50.59,279.52) and (53.99,278.27) .. (58.02,278.25) .. controls (62.98,278.22) and (67.01,280.07) .. (67.02,282.38) .. controls (67.03,283.36) and (66.32,284.25) .. (65.13,284.97) ;  
\draw  [draw opacity=0] (80.93,286.33) .. controls (78.89,285.48) and (77.6,284.22) .. (77.6,282.8) .. controls (77.6,280.26) and (81.76,278.2) .. (86.9,278.2) .. controls (90.96,278.2) and (94.41,279.49) .. (95.68,281.28) -- (86.9,282.8) -- cycle ; \draw   (80.93,286.33) .. controls (78.89,285.48) and (77.6,284.22) .. (77.6,282.8) .. controls (77.6,280.26) and (81.76,278.2) .. (86.9,278.2) .. controls (90.96,278.2) and (94.41,279.49) .. (95.68,281.28) ;  

\draw (16.4,180) node [anchor=north west][inner sep=0.75pt]  [font=\tiny] [align=left] {$\displaystyle \times $};
\draw (115.2,190) node [anchor=north west][inner sep=0.75pt]  [font=\tiny] [align=left] {$\displaystyle \times $};
\draw (64.4,143.2) node [anchor=north west][inner sep=0.75pt]  [font=\tiny] [align=left] {$\displaystyle \times $};
\draw (116.4,150.4) node [anchor=north west][inner sep=0.75pt]  [font=\tiny] [align=left] {$\displaystyle \times $};
\draw (59.6,207.2) node [anchor=north west][inner sep=0.75pt]  [font=\tiny] [align=left] {$\displaystyle \times $};
\draw (143.6,158.4) node [anchor=north west][inner sep=0.75pt]  [font=\tiny] [align=left] {$\displaystyle \times $};

\end{tikzpicture}

    \caption{Primer obešene slike. Podana shema ne reši problema.}
    \label{mucka}
\end{figure}

Tukaj omenimo, da je to najpreprostejša oblika problema predstavljenega v \cite{picture-hanging-puzzles}. Tega bi lahko posplošili sledeče: denimo, da imamo v steni zabitih $n$ žebljev, tedaj zahtevamo, da slika pade natanko tedaj, ko iz nje odstranimo točno $k$ žebljev. Seveda pa se, ker naš problem uporabimo le v ilustrativne namene, tukaj držimo enostavnejše verzije.

Kako pa se reševanja takega problema sploh lotiti?

V tem prispevku to naredimo preko fundamentalne grupe $n$-krat preluknjane evklidske ravnine. Razlog za to je, da je fundamentalna grupa izvrstno orodje za opisovanje globalne strukture topoloških prostorov. To doseže s pomočjo zank v prostoru in ekvivalenčno relacijo homotopnosti med njimi, ki nam pove, na katerih elementih prostora se te "{}zataknejo". Lahko si predstavljamo, da je taka struktura odlična za opisovanje lukenj v dotičnem prostoru.

Za nas bo to prišlo v poštev pri matematičnem modelu problema, saj, če si našo steno zamislimo kot evklidsko ravnino in vsako vrzel, ki smo jo v njej ustvarili, kot enega izmed žebljev, naš problem opiše ravno to, kje se zanke — torej različna navijanja vrvi okoli žebljev — zataknejo.

Za matematično analizo privzamemo naslednje splošne oznake: z $\N$ označujemo množico naravnih, z $\Z$ množico celih in z $\R$ množico realnih števil, z $\I$ označujemo enotski interval, z $\B^n$ odprto $n$-kroglo, s $\S^n$ $n$-sfero in s $\T^n$ $n$-torus. Poleg tega s $\F_n$ označimo prosto grupo ranga $n$.

\section{Poti v prostoru}


Za začetek se spomnimo nekaj klasičnih definicij s področja splošne topologije, ki jih bomo potrebovali v nadaljevanju.


\begin{definicija}
    Naj bo $X$ topološki prostor. Zvezna preslikava $\alpha : \I \to X$ je \textit{pot} v $X$. Element $x = \alpha(0)$ imenujemo \textit{začetna točka}, element $y = \alpha(1)$ pa \textit{končna točka} poti $\alpha$. Tedaj pravimo, da je $\alpha$ pot od $x$ do $y$.

    Poti $\alpha$ v $X$ s staknjenima krajiščema pravimo \textit{zanka v točki} $x = \alpha(0) = \alpha(1)$. Z $Z(X, x)$ označimo \textit{množico vseh zank v} $x$.
\end{definicija}


V prihodnje se ukvarjamo predvsem z zankami, kljub temu pa vseeno večino trditev dokažemo v splošnem, saj nam to dokazov ne oteži.

Nadaljujemo z nekaj operacijami.


\begin{definicija}
    Naj bosta $\alpha$ in $\beta$ poti v $X$, za kateri velja $\alpha(1) = \beta(0)$. Z $\alpha\beta$ označimo preslikavo $\I \to X$ podano s predpisom
    \begin{align*}
        x \mapsto
        \begin{cases}
            \alpha(2x) & (x \in [0, 1/2]) \\
            \beta(2x - 1) & (x \in [1/2, 1])
        \end{cases},
    \end{align*}
    ki ji pravimo \textit{stik poti} $\alpha$ in $\beta$.
\end{definicija}


Tukaj je vredno omeniti, da je stik zank spet zanka, saj je to pomembno v nadaljevanju.

Poleg tega pa potrebujemo tudi sledečo involucijo.


\begin{definicija}
    Naj bo $\alpha$ pot v $X$. Z $\alpha^{-1}$ označimo \textit{obratno pot} poti $\alpha$ definirano s preslikavo $\I \to X, x \mapsto \alpha(1 - x)$.
\end{definicija}


\begin{opomba}
    Obratno pot se pogosto označuje $\overline\alpha$, vendar jo mi označujemo kot jo v duhu definicij in izrekov v sledečih poglavjih.
\end{opomba}


Ker konstantne poti igrajo pomembno vlogo v izpeljavi, jih označujemo kar s točko, v katero slikajo. Torej naj bo $x \in X$, z $x$ označujemo preslikavo $\I \to X, s \mapsto x$.


\section{Homotopija}


Homotopija je eno najpomembnejših orodij algebraične topologije. Poleg tega, da igra pomembno vlogo pri konstrukciji homotopskih in homoloških grup, lahko z njo opišemo tudi mnoge druge homotopske invariante. Tukaj je ne razvijemo v vsej njeni kompleksnosti, saj za konstrukcijo fundamentalne grupe potrebujemo le najosnovnejše definicije.


\begin{definicija}
    Naj bosta $X$ in $Y$ topološka prostora in $\phi, \psi : X \to Y$ zvezni. Preslikavi $\phi$ in $\psi$ sta \textit{homotopni} natanko tedaj, ko obstaja zvezna preslikava $h : X \times \I \to Y$, za katero velja, da je $h(x, 0) = \phi(x)$ in $h(x, 1) = \psi(x)$ za vse $x \in X$. Tedaj pišemo $\phi \simeq \psi$, preslikavi $h$ pa pravimo \textit{homotopija} med $\phi$ in $\psi$.

    Naj bo $A \subseteq X$. Če obstaja kakšna homotopija $h$, za katero velja $h(x, y) = \phi(x) = \psi(x)$ za vse $x \in A$ in $y \in \I$, pravimo, da sta $\phi$ in $\psi$ \textit{homotopni relativno na} $A$. Pišemo $\phi \simeq \psi$ $(\rel A)$.
\end{definicija}


Homotopijo si najlažje predstavljamo kot zvezno transformacijo slike ene preslikave v drugo skozi topološki prostor. To pomeni, da ni važno le, kako vsaka slika izgleda, temveč tudi, kje se nahaja. Kot primer si poglejmo kanonični inkluziji $\phi, \psi : \B^n \to \B^n \amalg \B^n$, kjer $\phi$ slika v prvo, $\psi$ pa v drugo kopijo $\B^n$. Kljub temu, da se preslikavi zdita skoraj popolnoma enaki, pa še vseeno nista homotopni, saj nobene izmed žog ne moremo potisniti v drugo brez nenadnega skoka.

Kot si za relacije zelo pogosto želimo, je tudi homotopnost ekvivalenčna.


\begin{trditev}
    Homotopnost je ekvivalenčna relacija.
    \label{homotopnost-ekvivalencna-relacija}
\end{trditev}


\begin{proof}
    Naj bodo $\phi, \psi, \sigma : X \to Y$ zvezne. Očitno velja $\phi \simeq \phi$. Naj bo $h$ homotopija med $\phi$ in $\psi$. Tedaj je preslikava $(x, y) \mapsto h(x, 1 - y)$ homotopija med $\psi$ in $\phi$. Nazadnje naj bo $h$ kot prej in $h'$ homotopija med $\psi$ in $\sigma$. Preslikava definirana s
    \begin{align*}
        (x, y) \mapsto
        \begin{cases}
            h(x, 2y) & (y \in [0, 1/2]) \\
            h'(x, 2y - 1) & (y \in [1/2, 1])
        \end{cases}
    \end{align*}
    tvori homotopijo med $\phi$ in $\sigma$.
\end{proof}


Saj se želimo ukvarjati predvsem s potmi, moramo definicijo homotopnosti rahlo prilagoditi. Problem z zdajšno namreč nastane v s potmi povezanih prostorih, saj so v njih vse poti homotopne. Recimo, da je $X$ povezan s potmi in $\alpha, \beta$ poti v $X$. Ker obstaja neka pot $\sigma$ od $\alpha(1)$ do $\beta(0)$, lahko homotopijo med njima konstuiramo, tako da pot $\alpha$ po tiru $\sigma$ preprosto povlečemo na mesto poti $\beta$.

Opisanemu problemu se zato izognemo s sledečo definicijo.


\begin{definicija}
    Naj bosta $x, y \in X$ in $\alpha, \beta$ poti od $x$ do $y$. Poti $\alpha$ in $\beta$ sta \textit{homotopni glede na krajišča} natanko tedaj, ko sta homotopni relativno na $\{0, 1\}$.
\end{definicija}


\begin{opomba}
    Od tu naprej vedno, ko govorimo o homotopnosti poti, mislimo homotopnost glede na krajišča le-teh. Kljub temu pa je pomembno poudariti, da se moramo tega na vsakem koraku zavedati, saj velika večina izrekov in trditev temelji na tej predpostavki.
\end{opomba}


S tem, da se omejimo na homotopije, ki krajišča poti ohranjajo primrznjena, tako ustvarimo relacijo z dokaj drugačno strukturo.


\begin{zgled}
    Poglejmo si tri poti v podprostoru ravnine, kot kaže slika \ref{homotopija-poti}. Ni se težko prepričati, da sta poti $\alpha$ in $\beta$ homotopni, saj si lahko predstavljamo, kako bi ena zdrsela v drugo, kot bi jo s prstom potisnili v želeni položaj. Po drugi strani pa nobena izmed $\alpha$ in $\beta$ ni homotopna $\gamma$. Razlog je luknja v sredini prostora, ki preprečuje katerikoli od njiju zlesti na mesto poti $\gamma$. Da bi to dosegli, bi ju morali ali pretrgati ali pa v procesu transformacije začasno premakniti eno od krajišč, oba ta primera pa sta po definiciji homotopnosti poti nemogoča.

    \begin{figure}[H]
        \centering
        

\tikzset{every picture/.style={line width=0.75pt}}
\begin{tikzpicture}[x=1pt,y=1pt,yscale=-1,xscale=1]

\draw  [color={rgb, 255:red, 155; green, 155; blue, 155 }  ,draw opacity=1 ][fill={rgb, 255:red, 155; green, 155; blue, 155 }  ,fill opacity=0.5 ] (64.2,172.2) .. controls (103.4,158.2) and (189.4,156.6) .. (219,202.6) .. controls (248.6,248.6) and (161,219.4) .. (129.4,247) .. controls (97.8,274.6) and (73.8,318.2) .. (24.6,280.2) .. controls (-24.6,242.2) and (25,186.2) .. (64.2,172.2) -- cycle ;
\draw    (21.8,239) .. controls (28.6,217.8) and (71.4,157.8) .. (110.2,181.4) .. controls (149,205) and (183.8,179.8) .. (204.6,206.2) ;
\draw [color={rgb, 255:red, 0; green, 0; blue, 0 }  ,draw opacity=1 ]   (21.8,239) .. controls (56.6,236.6) and (52.2,207.4) .. (79.4,195.8) .. controls (106.6,184.2) and (106.6,217.4) .. (132.2,222.2) .. controls (157.8,227) and (195.18,213.27) .. (204.6,206.2) ;
\draw    (21.8,239) .. controls (21.8,252.2) and (18.2,264.6) .. (62.2,277.8) .. controls (106.2,291) and (99,232.2) .. (145.8,229.8) .. controls (192.6,227.4) and (206.6,221.8) .. (204.6,206.2) ;
\draw  [color={rgb, 255:red, 155; green, 155; blue, 155 }  ,draw opacity=1 ][fill={rgb, 255:red, 255; green, 255; blue, 255 }  ,fill opacity=1 ] (54.48,252.01) .. controls (49.28,243.68) and (57.23,229.33) .. (72.23,219.96) .. controls (87.24,210.59) and (103.62,209.75) .. (108.82,218.09) .. controls (114.02,226.42) and (106.08,240.77) .. (91.07,250.14) .. controls (76.07,259.51) and (59.69,260.35) .. (54.48,252.01) -- cycle ;
\draw  [fill={rgb, 255:red, 0; green, 0; blue, 0 }  ,fill opacity=1 ] (20.1,239) .. controls (20.1,238.06) and (20.86,237.3) .. (21.8,237.3) .. controls (22.74,237.3) and (23.5,238.06) .. (23.5,239) .. controls (23.5,239.94) and (22.74,240.7) .. (21.8,240.7) .. controls (20.86,240.7) and (20.1,239.94) .. (20.1,239) -- cycle ;
\draw  [fill={rgb, 255:red, 0; green, 0; blue, 0 }  ,fill opacity=1 ] (202.9,206.2) .. controls (202.9,205.26) and (203.66,204.5) .. (204.6,204.5) .. controls (205.54,204.5) and (206.3,205.26) .. (206.3,206.2) .. controls (206.3,207.14) and (205.54,207.9) .. (204.6,207.9) .. controls (203.66,207.9) and (202.9,207.14) .. (202.9,206.2) -- cycle ;

\draw (150,179.2) node [anchor=north west][inner sep=0.75pt][font=\footnotesize][align=left] {$\displaystyle \alpha $};
\draw (44.8,209.2) node [anchor=north west][inner sep=0.75pt][font=\footnotesize][align=left] {$\displaystyle \beta $};
\draw (78.8,262.4) node [anchor=north west][inner sep=0.75pt][font=\footnotesize][align=left] {$\displaystyle \gamma $};

\end{tikzpicture}

        \caption{Primer homotopnih in nehomotopnih poti v podprostoru evklidske ravnine.}
        \label{homotopija-poti}
    \end{figure}

    Podobno sta si zanki $\alpha\gamma^{-1}$ in $\beta\gamma^{-1}$ homotopni, medtem ko si zanki $\alpha\beta^{-1}$ in $\alpha\gamma^{-1}$ nista. Razlog spet leži v luknji sredi prostora, saj le-ta zanki $\alpha\gamma^{-1}$ preprečuje, da bi se lahko zategnila v konstantno pot, zanka $\alpha\beta^{-1}$ pa nima takšne omejitve.
    \label{homotopija-poti-v-podprostoru-R2}
\end{zgled}


\section{Fundamentalna grupa}


Sedaj, ko smo se seznanili z osnovami homotopije poti, pa nadaljujemo h konstrukciji fundamentalne grupe. To tvori prostor ekvivalenčnih razredov zank glede na relacijo homotopnosti opremljen z operacijo stika poti, zato se osredotočimo na nekaj lem, ki to konstrukcijo legitimirajo. Za začetek dokažemo, da je stik poti usklajen s homotopnostjo, kar nam dovoli, da ga podaljšamo tudi na ekvivalenčne razrede.


\begin{lema}
    Naj bodo $\alpha, \beta, \gamma, \delta$ poti v $X$, za katere velja $\alpha \simeq \gamma$ in $\beta \simeq \delta$. Če $\alpha\beta$ in $\gamma\delta$ obstajata, potem sta homotopna.
    \label{stik-poti-homotopen}
\end{lema}


\begin{proof}
    Ker sta $\alpha$ in $\gamma$ homotopni, se ni težko prepričati, da sta tudi $\alpha\beta$ in $\gamma\beta$, saj jima v resnici dodamo le primrznjen rep, ki na zveznost podaljšane homotopije ne more vplivati. Na enak način se prepričamo, da isto velja za $\gamma\beta$ in $\gamma\delta$. Po trditvi \ref{homotopnost-ekvivalencna-relacija} je homotopnost poti tranzitivna.
\end{proof}


Nadaljujemo z osnovnimi lastnostmi grupe.


\begin{lema} Naj bo $x \in X$ in $\alpha, \beta, \gamma$ zanke v $x$.
    \begin{enumerate}[(i)]
        \item \label{lastnosti-grupe-1} $(\alpha\beta)\gamma \simeq \alpha(\beta\gamma)$
        \item \label{lastnosti-grupe-2} $\alpha x \simeq x \alpha \simeq \alpha$,
        \item \label{lastnosti-grupe-3} $\alpha^{-1}\alpha \simeq \alpha\alpha^{-1} \simeq x$.
    \end{enumerate}
    \label{lastnosti-grupe}
\end{lema}


\begin{proof}
    Zaradi njune rutinske narave dokaza za \eqref{lastnosti-grupe-1} in \eqref{lastnosti-grupe-2} izpustimo. Za \eqref{lastnosti-grupe-3} naj bo $\alpha$ zanka v $X$. Brez škode splošnosti se lahko omejimo na primer $\alpha\alpha^{-1} \simeq x$. Če si zamislimo, kako $\alpha\alpha^{-1}$ potuje v $X$, hitro opazimo, da kljub temu, da se točke $x$ dotakne vsaj trikrat, je tam fiksirana samo za $0$ in $1$. Tako lahko tisto nevidno zankico, ki nastane pri $1 / 2$, samo vlečemo nazaj po tiru $\alpha$, dokler se ne vrnemo do točke $x$. V skladu z razmislekom definiramo preslikavo $h : \I^2 \to X, (s, t) \mapsto \alpha(t - t|2s - 1|)$, ki tvori želeno homotopijo.
\end{proof}


Vse te leme in definicije pa nas le pripeljejo do velikega izreka (relativno).


\begin{izrek}
    Naj bo $X$ topološki prostor in $x \in X$. Množica $Z(X, x) / \!\!\simeq$ opremljena z notranjo operacijo $([\alpha], [\beta]) \mapsto [\alpha\beta]$ tvori grupo.
\end{izrek}


\begin{proof}
    Sledi trivialno po lemah \ref{stik-poti-homotopen} in \ref{lastnosti-grupe}.
\end{proof}


In tako lahko končno definiramo.


\begin{definicija}
    Naj bo $X$ topološki prostor in $x \in X$. Grupi $Z(X, x) / \!\!\simeq$ pravimo \textit{prva homotopska} ali \textit{fundamentalna grupa} topološkega prostora $X$ in jo označimo $\pi_1(X, x)$.
\end{definicija}


Preden se lotimo modeliranja, bi si bilo dobro dobiti še malo občutka za pravkar definirano strukturo. Pa si poglejmo še nekaj primerov fundamentalnih grup, ki pripadajo klasičnim mnogoterostim.


\begin{zgled}
    Fundamentalna grupa $\S^2$ je trivialna. Predstavljajmo si eno tistih žogic iz gumic. Gumico na površini lahko vedno vzamemo dol iz žogice, pri čemer je naprej napeta, potem pa se v procesu odstranjevanja počasi sprošča ob površini žoge, dokler se na neki točki popolnoma ne sprosti in sama pade dol s kugle.

    Bolj konkretno naj bo $\alpha$ neka zanka v $x \in \S^2$ in predpostavimo, da obstaja neka točka $y \in \S^2 \setminus \im \alpha$. Predstavljajmo si, da sfero predremo tj. odstranimo $y$. Dobro znano je, da je $\S^2 \setminus \{y\} \approx \B^2$, naj bo $\phi : \S^2 \setminus \{y\} \to \B^2$ ta homeomorfizem. Predpostavimo lahko, da je $\phi(x) = 0$. Poleg tega je preslikava $H : \B^2 \times \I \to \B^2, (v, s) \mapsto (1 - s) v$ homotopija med $\id_{\B^2}$ in $\B^2 \mapsto \B^2, v \mapsto 0$. Sedaj lahko definiramo preslikavo $h : \I^2 \to \S^2$ s predpisom $(s, t) \mapsto \phi^{-1}(H((\phi \circ \alpha)(s), t))$, ki tvori homotopijo med $\alpha$ in $x$.

    Denimo, da tak $y$ ne obstaja tj. $\im \alpha = \S^2$. Da bi lahko uporabili predhonji argument, moramo naprej najti neko pot homotopno $\alpha$, ki nam bo to dovolila. Izberimo si zaprto $U \subset \S^2$ homeomorfno disku, za katero velja $x \notin U$. Poglejmo si, kaj naredimo z $\alpha$, ko seka $U$. Naj bo $I_\lambda = [u, v] \subset \I$ tak, da velja $\alpha(u), \alpha(v) \in \partial U$ in $\im \alpha|_{[u, v]} \subseteq U$, in $\partial_\lambda : \I \to \partial U$ pot, za katero velja $\partial_\lambda(0) = \alpha(u)$ in $\partial_\lambda(1) = \alpha(v)$. Poleg tega s $\phi_\lambda$ označimo nek homeomorfizem $\I \to I_\lambda$, za katerega velja $\phi_\lambda(0) = u$ in $\phi_\lambda(1) = v$. Ni se težko prepričati, da je $\alpha|_{[u, v]} \circ \phi_\lambda \simeq \partial_\lambda$ v $U$. Homotopijo med njima označimo s $H_\lambda$.

    Naj bo $\{I_\lambda\}_{\lambda \in \Lambda}$ družina vseh maksimalnih intervalov, kot so opisani zgoraj, $\{\partial_\lambda\}_{\lambda \in \Lambda}$ njim pripadajoče poti in $\{\phi_\lambda\}_{\lambda \in \Lambda}$ homeomorfizmi. Pot $\sigma : \I \to \S^2$ definiramo s predpisom $\sigma(t) = \partial_\lambda(t)$, ko je $t \in I_\lambda$, in $\sigma(t) = \alpha(t)$ sicer. Preslikava $h : \I^2 \to \S^2$, ki jo definiramo s predpisom $h(s, t) = H_\lambda(\phi_\lambda^{-1}(s), t)$ za $s \in I_\lambda$ in $h(s, t) = \alpha(s)$ sicer, je homotopija med $\alpha$ in $\sigma$.

    Pot $\sigma$ je bila konstruirana tako, da smo $\alpha$ preprosto odstranili iz $\interior U$. Ker pa $\interior U \neq \emptyset$, si lahko izberemo nek $y \notin \im\sigma$, zato je $\sigma \simeq x$, posledično pa velja tudi $\alpha \simeq x$.
    \label{fundamentalna-grupa-S2}
\end{zgled}


Pri prejšnjem zgledu je vredno omeniti, da je celotno konstrukcijo ostala neodvisna od izbire bazne točke, posledično pa je potem taka tudi $\pi_1(\S^2)$. Ta lastnost pa ni edinstvena $\S^2$, temveč velja za vse s potmi povezane topološke prostore.


\begin{trditev}
    Naj bo $X$ povezan s potmi in $x, y \in X$. Tedaj velja $\pi_1(X, x) \cong \pi_1(X, y)$.
\end{trditev}


\begin{proof}
    Naj bo $\sigma$ pot od $x$ do $y$ v $X$. Poglejmo si preslikavo $\phi : \pi_1(X, x) \to \pi_1(X, y), [\alpha] \mapsto [\sigma^{-1}\alpha\sigma]$. Ni težko preveriti, da je $\phi$ homomorfizem. Tako nam ostane le še bijektivnost. Definirajmo $\psi : \pi_1(X, y) \to \pi_1(X, x), [\alpha] \mapsto [\sigma\alpha\sigma^{-1}]$. Velja $\psi \circ \phi = \id$ in $\phi \circ \psi = \id$, torej je $\phi$ bijekcija.
\end{proof}


\begin{opomba}
    V primeru, da ni odvisna od izbire bazne točke, fundamentalno grupo prostora $X$ označujemo kar $\pi_1(X, x)$. To ne pomeni nujno, da je $X$ povezan s potmi. Kot primer si lahko vzamemo $\pi_1(\B^n \coprod \B^n)$.
\end{opomba}


\begin{zgled}
    Poglejmo si prostore $\S^1$, $A = \S^1 \times \I$ in $K = \S^1 \times \B^2$. Začnemo s $\S^1$. Naj bo $\alpha : \I \to \S^2$ definirana s predpisom $t \mapsto e^{i2\pi t}$. Hitro opazimo, da velja $\im\alpha = \S^1$. Tako se moramo pri ustvarjanju $\alpha$ homotopnih poti omejiti na premikanje po tiru $\alpha$ tj. $\alpha$ lahko vzdolž njega le raztegujemo in vlečemo. Sledi, da $\alpha$ ne more biti homotopna konstanti, saj se je omejene na svojo tir nikakor ne da pripraviti, da bi preskočila luknjo v $\S^1$. Podobno $\alpha$ ne more biti homotopna $\alpha\alpha$ ali $\alpha^{-1}$. S tem v mislih je jasno, da morajo biti vsi elementi $\pi_1(\S^1)$ oblike $[\alpha^n]$ za nek $n \in \Z$. Od tu naravno sledi $\pi_1(\S^1) \cong \Z$.

    Podobno kot pri $\S^1$ sta tudi $A$ in $K$ generirana z enim elementom. Enako kot prej lahko definiramo $\alpha : \I \to A, \beta : \I \to K$ s predpisom $t \mapsto (e^{i2\pi t}, 0)$, nakar nam bosta $[\alpha]$ in $[\beta]$ služila kot generatorja vsak svoje grupe. Ni pretirano težko videti, da morata biti tudi edina. Kot primer si oglejmo neko poljubno $\gamma = (\gamma_1, \gamma_2)$ v $K$. Naj bo $H$ definirana kot v prvem delu zgleda \ref{fundamentalna-grupa-S2}. Tedaj je $h : \I^2 \to K, (x, y) \mapsto (\gamma_1(x), H(\gamma_2(x), y))$ homotopija med $\gamma$ in $(\gamma_1, 0)$, ki leži na tiru $\beta$, potem smo pa spet na istem, kot če bi imeli le $\S^1$. To pomeni, da velja $\pi_1(\S^1) \cong \pi_1(A) \cong \pi_1(K)$.
    \label{fundamentalna-grupa-S1}
\end{zgled}


\begin{zgled}
    Zanima nas, kaj je fundamentalna grupa mnogoterosti $\T^2$. Izgleda, kot da ima $\pi_1(\T^2)$ natanko dva generatorja — zanko, ki gre po obodu luknje, in zanko, ki obkroži biskvit, kot kaže slika \ref{krof} —, poleg tega noben izmed njiju ni končnega reda, torej bo $\pi_1(\T^2)$ sigurno neskončna. Preverimo, ali je komutativna tj. ali je izomorfna $\Z^2$.

    \begin{figure}[H]
        \centering
        

\tikzset{every picture/.style={line width=0.75pt}}
\begin{tikzpicture}[x=1.5,y=1.5pt,yscale=-1,xscale=1]

\draw  [color={rgb, 255:red, 155; green, 155; blue, 155 }  ,draw opacity=1 ][fill={rgb, 255:red, 155; green, 155; blue, 155 }  ,fill opacity=0.5 ][line width=0.75]  (27.29,219) .. controls (47.49,204.2) and (114.51,201.91) .. (126.71,218.71) .. controls (138.91,235.51) and (140.45,247.05) .. (121.25,260.25) .. controls (102.05,273.45) and (48.85,280.25) .. (34.25,270.25) .. controls (19.65,260.25) and (7.09,233.8) .. (27.29,219) -- cycle ;
\draw  [color={rgb, 255:red, 155; green, 155; blue, 155 }  ,draw opacity=1 ][fill={rgb, 255:red, 255; green, 255; blue, 255 }  ,fill opacity=1 ][line width=0.75]  (56.25,239.54) .. controls (56.25,239.54) and (63.83,244.37) .. (73.42,243.46) .. controls (83,242.54) and (90.17,234.37) .. (90.17,234.37) .. controls (90.17,234.37) and (84.02,227) .. (71.86,229) .. controls (59.69,231) and (56.25,239.54) .. (56.25,239.54) -- cycle ;
\draw [color={rgb, 255:red, 155; green, 155; blue, 155 }  ,draw opacity=1 ][line width=0.75]    (90.17,234.37) .. controls (92,232.46) and (92.75,231.29) .. (92.75,231.29) ;
\draw [color={rgb, 255:red, 155; green, 155; blue, 155 }  ,draw opacity=1 ][line width=0.75]    (52.25,235.79) .. controls (52.25,235.79) and (52,236.96) .. (56.25,239.54) ;
\draw   (78.43,257.58) .. controls (76.11,248.83) and (79.43,240.36) .. (85.84,238.66) .. controls (92.25,236.96) and (99.33,242.68) .. (101.64,251.43) .. controls (103.96,260.18) and (100.65,268.65) .. (94.24,270.35) .. controls (87.83,272.05) and (80.75,266.33) .. (78.43,257.58) -- cycle ;
\draw  [draw opacity=0] (91.78,270.66) .. controls (86.08,270.59) and (80.4,265.21) .. (78.41,257.51) .. controls (76.39,249.67) and (78.87,242.08) .. (84.01,239.38) -- (90.04,254.5) -- cycle ; \draw  [densely dashed, color={rgb, 255:red, 255; green, 255; blue, 255 }  ,draw opacity=1 ] (91.78,270.66) .. controls (86.08,270.59) and (80.4,265.21) .. (78.41,257.51) .. controls (76.39,249.67) and (78.87,242.08) .. (84.01,239.38) ;  
    \draw  [draw opacity=0] (91.78,270.66) .. controls (86.08,270.59) and (80.4,265.21) .. (78.41,257.51) .. controls (76.39,249.67) and (78.87,242.08) .. (84.01,239.38) -- (90.04,254.5) -- cycle ; \draw  [densely dashed, color={rgb, 255:red, 155; green, 155; blue, 155 }  ,draw opacity=0.5 ] (91.78,270.66) .. controls (86.08,270.59) and (80.4,265.21) .. (78.41,257.51) .. controls (76.39,249.67) and (78.87,242.08) .. (84.01,239.38) ;  
\draw   (30.46,240.32) .. controls (29.06,228.73) and (47.54,216.95) .. (71.74,214.01) .. controls (95.94,211.07) and (116.7,218.09) .. (118.11,229.68) .. controls (119.51,241.27) and (101.03,253.05) .. (76.83,255.99) .. controls (52.63,258.93) and (31.87,251.91) .. (30.46,240.32) -- cycle ;

\draw (103,257) node [anchor=north west][inner sep=0.75pt][font=\footnotesize][align=left] {$\displaystyle \alpha $};
\draw (30,250) node [anchor=north west][inner sep=0.75pt][font=\footnotesize][align=left] {$\displaystyle \beta $};

\end{tikzpicture}

        \caption{Generatorja fundamentalne grupe prostora $\T^2$.}
        \label{krof}
    \end{figure}

    Naj bosta $\alpha, \beta$ generatorja $\pi_1(\T^2)$. Če si torus predstavljamo kot kvocient $\I^2$, je lahko videti, da velja $\alpha\beta \simeq \beta\alpha$, saj, če dobljeno zanko primemo, ko se ponovno vrne v $x$, in jo povlečemo skozi kot $\I^2$, se zaradi ekvivalenčne relacija spet pojavi v nasprotnem. Iz tega sledi, da je fundamentalna grupa torusa komutativna, od tod pa lahko zaključimo, da velja $\pi_1(\T^2) \cong \Z^2$.
\end{zgled}


Množico $\T^n$ se pogosto definira kar kot produkt $(\S^1)^n$. Pravkar smo pokazali, da velja $\pi_1(\S^1 \times \S^1) = \Z^2$, vendar, ker velja $\pi_1(\S^1) = \Z$, sledi $\pi_1(\S^1 \times \S^1) \cong \pi_1(\S^1) \times \pi_1(\S^1)$. Poleg tega pa, če se vrnemo k zgledu \ref{fundamentalna-grupa-S1}, opazimo, da spet velja $\pi_1(\S^1 \times \I) \cong \pi_1(\S^1) \times \pi_1(\I) \cong \pi_1(\S^1)$. Isti vzorec opazimo tudi pri $\pi_1(K)$. V vseh teh primerih velja, da je fundamentalna grupa produkta enaka produktu fundamentalnih grup. Izkaže se, da to velja tudi v splošnem.


\begin{trditev}
    Naj bosta $X, Y$ topološka prostora, $x \in X$ in $y \in Y$. Tedaj velja $\pi_1(X \times Y, (x, y)) \cong \pi_1(X, x) \times \pi_1(Y, y)$.
\end{trditev}


\begin{proof}
    Definiramo preslikavo $\phi : \pi_1(X \times Y, (x, y)) \to \pi_1(X, x) \times \pi_1(Y, y)$ s predpisom $[\alpha] \mapsto ([\pi_X \circ \alpha], [\pi_Y \circ \alpha])$. Dovolj je pokazati, da je $\phi$ injektivna in da sta obe komponenti $\phi = (\phi_X, \phi_Y)$ surjektivna homomorfizma. Brez škode splošnosti se fokusiramo le na $\phi_X$. Izberimo si nek $[\alpha] \in \pi_1(X, x)$, velja $\phi_X([(\alpha, y)]) = [\alpha]$. Poleg tega pa naj bosta tokrat $[\alpha], [\beta] \in \pi_1(X \times Y, (x, y))$. Velja $\phi_X([\alpha\beta]) = [\pi_X \circ \alpha\beta] = [\pi_X \circ \alpha][\pi_X \circ \beta] = \phi_X([\alpha])\phi_X([\beta])$, od koder sledi, da je $\phi_X$ homomorfizem. Preostane nam le še injektivnost. Naj bosta $[\alpha]$ in $[\beta]$ kot prej. Razvijemo $\phi([\alpha]) = \phi([\beta])$ v $([\pi_X \circ \alpha], [\pi_Y \circ \alpha]) = ([\pi_X \circ \beta], [\pi_Y \circ \beta])$, kar jasno implicira $[\alpha] = [\beta]$.
\end{proof}


\section{Obesimo sliko}


Oboroženi z orodji algebraične topologije se le obrnemo proti modeliranju problema.

\begin{zgled}
    Kot že rečeno, si želimo za model stene uporabiti evklidsko ravnino. Naj bo $x \in \R^2$ in $S \subset \R^2 \setminus \{x\}$, kjer je $|S| = n$ za nek $n \in \N$. Izbrana točka $x$ nam bo služila kot mesto, na katerem je vrvica pritrjena na sliko, elementi $S$ predstavljajo vsak enega izmed žebljev zabitih v steno, za vrv pa bomo uporabili kar ekvivalenčni razred poti v prostoru glede na relacijo homotopnosti, kot že motivirano. S tako nastavljenim modelom je jasno, da bo slika padla, natanko takrat, ko bo naša vrv enaka ekvivalenčnemu razredu $[x]$, saj se takrat vrvica ni zataknila ob nobenega izmed žebljev. Torej našo steno opisuje grupa $G = \pi_1(\R^2 \setminus S, x)$.

    Grupa $G$ mora za vsak $s \in S$ vsebovati ekvivalenčni razred vseh zank, ki potujejo okoli $s$ natanko enkrat in ne obkrožijo nobenega drugega elementa. Taka ekvivalenčna razreda sta dva tj. naj bo $\alpha$ zanka kot zahtevano, potem sta $[\alpha]$ in $[\alpha^{-1}]$ iskana razreda. S preprostostjo v mislih ju označujemo kar s $s$ in $s^{-1}$. Vidimo, da je vsak element v $G$ možno napisati kot produkt teh osnovnih komponent. Sledi, da je $G = \{\prod_i a_i : a_i \in \bigcup_{s \in S} \{s, s^{-1} \} \}$. Intuitivno smo vsako možno navijanje vrvi zapisali s koraki, ki smo jih naredili, da smo ga ustvarili. Pokazali smo, da je možno $G$ zapisati kot množico vseh besed v $S$. Drugače rečeno velja $G \cong \F_n$.

    Preostane nam še opis odstranjevanja enega izmed žebljev. Naj bo $\alpha$ pot, ki potuje okoli točke $s \in S$ in nobene druge. Odstraniti žebelj pomeni odstraniti oviro — v našem primeru to pomeni zapolniti točko $s$, nakar bo $\alpha \simeq x$, torej bo slika padla. Ni nujno, da bo $\alpha$ vedno tako preprosta, po principu odstranjevanja ovire pa si še vedno želimo, da je po transformaciji prostora homotopna $x$.

    Pa to idejo formalizirajmo. Naj bo $\F_n = \langle S \rangle$. V skladu s povedanim za $s \in S$ definiramo preslikavo $\phi_s : \F_n \to \F_n$ s predpisom
    \begin{align*}
        t \mapsto
        \begin{cases}
            e & (t = s) \\
            t & (t \neq s)
        \end{cases},
    \end{align*}
    kjer je $e$ identiteta $\F_n$ in $t \in S$. Vsaka izmed preslikav $\phi_s$ opisuje odstranjevanje enega izmed žebljev. Seveda pa je v takem stanju definicija $\phi_s$ zelo pomanjkljiva, saj bi z njo znali operirati le na generatorij. Zato jo za preostale elemente $\F_n$ preprosto podaljšamo s predpisom $\phi_s(\prod_i a_i) = \prod_i \phi_s(a_i)$, kjer so $a_i \in S$. Preslikava $\phi_s$ je očitno homomorfizem grup.

    Preslikave $\phi_s$ so bile konstruirane na način, da se vsa navijanja, ki padejo, ko odstranimo žebelj $s$, nahajajo v podgrupi $\ker \phi_s$. Od tu sledi, da bodo rešitve ležale v preseku vseh, saj so to ravno vsa obešanja, ki padejo, ko odstranimo kateregakoli izmed žebljev. Torej je $\bigcap_{s \in S} \ker \phi_s$ množica rešitev problema obešanja slike.
    \label{konstrukcija-stene}
\end{zgled}


Na tej točki pa nam preostane le še dokazati, da te rešitve res obstajajo.


\begin{izrek}
    Naj bo $\F_n = \langle S \rangle$ in preslikava $\phi_s$ kot definirana v zgledu \ref{konstrukcija-stene} za vsak $s \in S$. Podgrupa $\bigcap_{s \in S} \ker \phi_s$ je netrivialna.
\end{izrek}


\begin{proof}
    Pričnimo s primerom za $n = 2$. Naj bosta $x, y$ generatorja $\F_2$. Označimo komutator $[x, y] = xyx^{-1}y^{-1}$. Ker $\F_n$ ni komutativna, velja $[x, y] \neq e$. Opazimo, da sta $\phi_x([x, y]), \phi_y([x, y]) = e$, torej je $\ker \phi_x \cap \ker \phi_y$ netrivialna.

   V splošnem naj bodo $x_i$ generatorji $\F_n$. Induktivno definiramo $S_1 = x_1$ in $S_p = [S_{p - 1}, x_p]$. Nadaljujemo z matematično indukcijo. Primer za $p = 1$ velja trivialno. Za indukcijski korak denimo, da je $S_p \in \ker \phi_i$ za vsak $i \leq p$. Tedaj je $\phi_{p + 1}(S_{p + 1}) = S_p S_p^{-1} = e$, saj $S_p$ ne vsebuje elementa $x_{p + 1}$. Poleg tega pa za vsak $i \leq p$ velja $\phi_i(S_{p + 1}) = x_{p + 1}x_{p + 1}^{-1} = e$. Torej je $S_{p + 1} \in \bigcap_{i \leq p + 1} \ker \phi_i$, kot želeno.

    Sledi, da je $S_n \in \bigcap_{s \in S} \ker \phi_s$ kot zahtevano.
\end{proof}


\printbibliography


\end{document}
